% lemniro: {"kind":"note","description":"A rendering reference for mathematical lecture notes: structured text, AMS mathematics, references, tables, and TikZ.","topic":"Topology","date":"2026-10-01","readingMinutes":5,"prerequisites":[],"related":[]}
\documentclass[11pt]{article}
\input{tex/lemniro-preamble}
\usepackage{tikz}
\usepackage{mathtools,cancel}
\usetikzlibrary{arrows.meta,positioning}
\newtheorem{claim}[theorem]{Claim}
\newcommand{\inv}[1]{#1^{-1}}
\title{Continuity of $f$: a complete lecture-note example}
\begin{document}
\maketitle

\section{Definitions and statements}
This document is ordinary LaTeX. Its paragraphs, statements, lists, and
references become HTML; a TikZ figure is compiled to vector graphics.
Text remains selectable and wraps to the width of the reader's screen.

\begin{definition}[Continuity]\label{def:continuous}
Let $X$ and $Y$ be topological spaces. A map $f:X\to Y$ is continuous if
$\inv{f}(V)$ is open in $X$ for every open subset $V\subseteq Y$.
\end{definition}

\begin{lemma}\label{lem:inverse}
For any function $f:X\to Y$ and any family $(V_\alpha)_{\alpha\in A}$,
\begin{equation}\label{eq:inverse}
  \inv{f}\!\left(\bigcup_{\alpha\in A}V_\alpha\right)
  =\bigcup_{\alpha\in A}\inv{f}(V_\alpha).
\end{equation}
\end{lemma}
\begin{proof}
A point $x$ lies in the left-hand side of~\eqref{eq:inverse} exactly when
$f(x)\in V_\alpha$ for some $\alpha\in A$. This holds exactly when
$x\in\inv{f}(V_\alpha)$ for some $\alpha\in A$, which means that $x$
lies in the right-hand side.
\end{proof}

\begin{theorem}\label{thm:composition}
The composition of two continuous maps is continuous.
\end{theorem}
\begin{proof}
Let $f:X\to Y$ and $g:Y\to Z$ be continuous, and let $W\subseteq Z$
be open. Then $\inv{g}(W)$ is open in $Y$ and
$\inv{f}(\inv{g}(W))$ is open in $X$. Since
\[
  \inv{(g\circ f)}(W)=\inv{f}\bigl(\inv{g}(W)\bigr),
\]
Definition~\ref{def:continuous} applies to $g\circ f$.
\end{proof}
\begin{corollary}
Every finite composition of continuous maps is continuous.
\end{corollary}
\begin{remark}
Theorem~\ref{thm:composition} does not require metrics on the spaces.
\end{remark}
\begin{claim}
The identity map of any topological space is continuous.
\end{claim}
\begin{example}
For a discrete domain, every function is continuous. For an indiscrete
codomain, every function is continuous as well.
\end{example}
\begin{exercise}
Use Definition~\ref{def:continuous} to prove the preceding example.
\end{exercise}

\section{A TikZ diagram}
The diagram in Figure~\ref{fig:composition} records the same composition.
\begin{figure}[ht]
\centering
\begin{tikzpicture}[>=Stealth, every node/.style={font=\large}]
  \node (X) at (0,0) {$X$};
  \node (Y) at (3,0) {$Y$};
  \node (Z) at (6,0) {$Z$};
  \draw[->,thick] (X) -- node[above] {$f$} (Y);
  \draw[->,thick] (Y) -- node[above] {$g$} (Z);
  \draw[->,thick] (X) to[bend right=30] node[below] {$g\circ f$} (Z);
\end{tikzpicture}
\caption{Continuous maps and their composition.}\label{fig:composition}
\end{figure}

\section{Mathematical typography}
The following identities exercise fractions, limits, roots, matrices,
cases, calligraphic letters, and blackboard bold alphabets.
\begin{align}
  \sum_{k=0}^{\infty}r^k &= \frac{1}{1-r}, && |r|<1,\label{eq:geometric}\\
  \int_0^1 x^2\,dx &= \frac13.\nonumber
\end{align}
\[
  A=\begin{pmatrix}1&2\\0&1\end{pmatrix},\qquad
  |x|=\begin{cases}x,&x\geq0,\\-x,&x<0,\end{cases}
  \qquad \sqrt{1+x^2}.
\]
In topology, $\T\subseteq\powerset{X}$, $\closure{U}\subseteq X$,
and $\Int(U)\subseteq U$. In algebra, $\mathbb{Q}\subseteq\R\subseteq\mathbb{C}$.
Equation~\eqref{eq:geometric} is a geometric-series identity.

\begin{example}[Macro equivalence]
The macro $\inv{f}$ and the literal expression $f^{-1}$ have identical
typography. Inline descenders $g_j$, fractions $\frac{a}{b}$, and script
letters $\mathcal{T}$ share the surrounding text's baseline.
\end{example}
\[
  X\xrightarrow{\text{continuous}}Y,
  \qquad \cancel{x}+x,
  \qquad \underbrace{x+\cdots+x}_{n\text{ terms}}.
\]
\begin{gather}
  \inv{f}(\varnothing)=\varnothing,\label{eq:empty}\\
  \inv{f}(Y)=X.\label{eq:whole}
\end{gather}
\begin{multline}\label{eq:long}
  \inv{f}(V_1\cup V_2\cup V_3\cup V_4)\\
  =\inv{f}(V_1)\cup\inv{f}(V_2)\cup\inv{f}(V_3)\cup\inv{f}(V_4).
\end{multline}
Equations~\eqref{eq:empty}, \eqref{eq:whole}, and~\eqref{eq:long}
retain their numbers and references.

\subsection{Lists, tables, and a footnote}
\begin{enumerate}
\item Start with a definition.
\item Prove a statement.
  \begin{itemize}\item Explain where each hypothesis is used.\end{itemize}
\item Work through an example.\footnote{Examples help distinguish a definition from its consequences.}
\end{enumerate}
\begin{center}
\begin{tabular}{ll}
\textbf{Construction} & \textbf{Open sets}\\
Discrete topology & Every subset\\
Indiscrete topology & $\varnothing$ and $X$
\end{tabular}
\end{center}
\end{document}
